\chapter{Língua Portuguesa}
\label{cap:portugues-teoria}

\begin{edital}
Compreensão e interpretação de textos; tipos e gêneros textuais; ortografia;
coesão (referenciação, conectores, tempos/modos verbais); estrutura
morfossintática (classes de palavras, coordenação, subordinação, pontuação,
concordância verbal e nominal, regência, crase, colocação pronominal);
reescrita e significação.
\end{edital}

Este capítulo cobre cada tópico do edital do zero, com a definição, o
raciocínio por trás da regra e exemplos resolvidos. A versão resumida ---
focada só no que costuma virar pegadinha de prova --- está no capítulo
correspondente da \texttt{apostila/}; aqui o objetivo é fixar o conceito de
verdade, não só reconhecer o padrão de banca.

\section{Compreensão e interpretação de textos}

\subsection{Três níveis de leitura}

Toda questão de interpretação testa um destes três níveis, e a maior parte
dos erros acontece por confundir um com o outro:

\begin{conceito}[Informação explícita, inferência e opinião]
\begin{itemize}
\item \textbf{Informação explícita}: está escrita literalmente no texto. A
pergunta se responde apontando a frase. Erro típico: não achar a frase por
estar em outro parágrafo, ou porque o examinador parafraseou (trocou as
palavras mas manteve o sentido).
\item \textbf{Inferência}: não está escrita, mas é a única conclusão possível
a partir do que está escrito + conhecimento de mundo compartilhado. Uma
inferência válida é \textbf{obrigatória}, não uma entre várias leituras
possíveis --- se dá pra imaginar um cenário em que a inferência é falsa, ela
não é uma inferência válida, é um chute.
\item \textbf{Opinião/juízo de valor}: pertence ao autor do texto (ou a um
personagem citado), não ao leitor nem ao examinador. A alternativa erra
quando atribui ao autor uma opinião que ele não expressou, ou quando o
candidato concorda/discorda pessoalmente do texto em vez de reportar o que o
texto diz.
\end{itemize}
\end{conceito}

\begin{exemplo}[Explícito vs.\ inferência vs.\ extrapolação]
Texto: \textit{"Depois de três tentativas frustradas de aprovar o projeto, a
diretoria decidiu arquivá-lo."}

\begin{itemize}
\item Explícito: "a diretoria arquivou o projeto" --- está escrito.
\item Inferência válida: "o projeto não foi aprovado nas tentativas
anteriores" --- não está escrito com essas palavras, mas é a única leitura
possível de "três tentativas frustradas".
\item Extrapolação (não é inferência válida): "a diretoria considera o
projeto inviável para sempre" --- o texto não diz isso; arquivar não é o
mesmo que descartar definitivamente. É o tipo de leitura que parece razoável
mas vai além do que o texto sustenta.
\end{itemize}
\end{exemplo}

\begin{cuidado}[Extrapolar não é inferir]
A diferença entre inferência válida e extrapolação é a diferença entre
"a única leitura possível" e "uma leitura plausível". Toda vez que uma
alternativa usa um verbo mais forte que o texto (o texto "sugere" e a
alternativa diz que ele "prova"; o texto descreve uma dificuldade e a
alternativa diz que é "impossível"), desconfie: é extrapolação.
\end{cuidado}

\subsection{Tipos textuais x gêneros textuais}

Esse par é cobrado quase sempre em conjunto e a diferença entre os dois
conceitos é a base de qualquer questão sobre "que tipo de texto é este":

\begin{conceito}[Tipo é a estrutura; gênero é a embalagem social]
\textbf{Tipo textual} é uma categoria \textbf{estrutural}, definida pela
natureza linguística predominante da composição --- são poucos e fixos:

\begin{itemize}
\item \textbf{Narração}: fatos organizados no tempo, com personagens, enredo,
conflito (verbos de ação, tempo passado predominante).
\item \textbf{Descrição}: apresenta características de algo/alguém, sem
progressão temporal (predomínio de adjetivos, verbos de estado).
\item \textbf{Dissertação/argumentação}: defende um ponto de vista com
argumentos, geralmente sobre um tema abstrato.
\item \textbf{Exposição}: explica, informa, expõe um assunto sem
necessariamente defender uma tese (diferença sutil com a dissertação: expor é
informar, dissertar é convencer).
\item \textbf{Injunção}: instrui, ordena, orienta uma ação (verbos no
imperativo ou infinitivo).
\end{itemize}

\textbf{Gênero textual} é a forma \textbf{social e concreta} que um texto
assume no mundo real, definida pela função comunicativa e pelo suporte ---
são inúmeros e mudam com o tempo (um gênero pode nascer com uma nova
tecnologia, como "post de rede social"): notícia, editorial, bula de remédio,
receita culinária, carta, e-mail, resenha, ata de reunião, edital,
regulamento, poema, conto, contrato.
\end{conceito}

\begin{exemplo}[Um gênero mistura tipos]
Uma \textbf{bula de remédio} (gênero) é predominantemente \textbf{injuntiva}
("tome 1 comprimido a cada 8 horas") mas tem trechos \textbf{descritivos}
("comprimido branco, redondo") e \textbf{expositivos} ("o princípio ativo age
bloqueando..."). Isso é a prova de que tipo e gênero são eixos diferentes: um
gênero (categoria fixa, social) é composto de vários tipos (categorias
estruturais) misturados.
\end{exemplo}

\begin{regrapratica}
Se a pergunta é "que \textbf{tipo}" de trecho é esse, procure a estrutura
predominante (tem verbo de ação e linha do tempo? narração. Só adjetivos,
sem ação? descrição). Se a pergunta é "que \textbf{gênero}" é o texto,
procure o nome social que ele teria fora da prova (editorial, e-mail, bula).
Uma reportagem de jornal (gênero) é majoritariamente narrativa com trechos
expositivos (tipo) --- a pergunta certa nunca mistura os dois níveis.
\end{regrapratica}

\subsection{Coerência x coesão}

\begin{conceito}[Dois requisitos diferentes de um texto bem construído]
\textbf{Coerência} é a lógica \textbf{do sentido}: as ideias do texto não se
contradizem e formam um todo que faz sentido para quem lê --- é uma
propriedade semântica/pragmática, depende do conhecimento de mundo do leitor.
\textbf{Coesão} é a amarração \textbf{gramatical} entre as partes do texto:
os recursos linguísticos visíveis (pronomes, conectores, repetição
controlada) que ligam frase a frase, parágrafo a parágrafo.

Um texto pode ser \textbf{coeso e incoerente}: todas as frases conectadas
gramaticalmente, mas o conteúdo é contraditório ou absurdo. E pode ser
\textbf{coerente e pouco coeso}: o sentido geral se entende, mas faltam
conectivos, e a leitura fica truncada. Os dois são independentes.
\end{conceito}

\section{Ortografia}

A ortografia da FGV cai pouco isoladamente (mais como parte de reescrita ou
concordância) porque o vigente é o Acordo Ortográfico de 1990, já
consolidado no ensino desde 2016. Os pontos que ainda geram dúvida real:

\begin{conceito}[Casos que sobrevivem como pegadinha]
\begin{itemize}
\item \textbf{Trema}: extinto para o português do Brasil (não se escreve mais
"lingüiça", "seqüência" --- ficou "linguiça", "sequência"). Só permanece em
nomes próprios e suas derivações (Müller, mülleriano).
\item \textbf{Hífen em palavras compostas com prefixo}: a regra prática (não
memorize todos os prefixos, olhe a \textbf{terminação do prefixo} x
\textbf{início da palavra seguinte}): mesma letra em ambos $\to$ hífen
("micro-ondas", "anti-inflamatório", com prefixo terminado em vogal + palavra
iniciada pela mesma vogal, ou prefixo terminado em consoante + palavra
iniciada pela mesma consoante, como "inter-regional"); prefixo terminado em
vogal + palavra iniciada em \textbf{r} ou \textbf{s} $\to$ dobra a letra, sem
hífen ("antirreligioso", "contrassenso"); demais casos, funde sem hífen
("autoescola", "extraoficial").
\item \textbf{Acentuação de ditongos abertos "ei" e "oi" em paroxítonas}: não
se acentuam mais ("ideia", "plateia", "heroico", "jiboia" --- antes eram
"idéia", "platéia", "heróico"). Continuam acentuados os ditongos abertos em
\textbf{oxítonas} ("herói", "chapéu", "céu").
\item \textbf{Acento diferencial}: caiu na reforma, com poucas exceções que
sobrevivem (\textbf{pôde} passado x \textbf{pode} presente; \textbf{pôr}
verbo x \textbf{por} preposição). "Para" (verbo parar, 3ª pessoa) e "pára"
não se distinguem mais por acento.
\end{itemize}
\end{conceito}

\begin{cuidado}[Reforma ortográfica não é "fonte volátil qualquer"]
O Acordo Ortográfico de 1990 está em vigor no Brasil desde 1º de janeiro de
2016 (fim do período de transição). Se uma questão datada de antes de 2016
usar "idéia" ou "vôo" com acento, não é erro do material --- é o texto
escrito sob a norma anterior. Não corrija exemplo histórico citado como
histórico.
\end{cuidado}

\section{Coesão textual: referenciação, conectores e tempos verbais}

\subsection{Coesão referencial}

\begin{conceito}[Anáfora, catáfora e elipse]
São os três mecanismos que evitam repetir a mesma palavra o tempo todo:

\begin{itemize}
\item \textbf{Anáfora}: um termo retoma algo já dito \textbf{antes}. É o caso
mais comum. \textit{"O relatório chegou atrasado. \textbf{Ele} continha
erros."} ("ele" retoma "o relatório").
\item \textbf{Catáfora}: um termo \textbf{anuncia} algo que só vem depois.
\textit{"Só quero uma coisa: \textbf{que você estude}."} (o "isso" implícito
aponta para frente).
\item \textbf{Elipse}: o termo é \textbf{omitido}, mas recuperável pelo
contexto. \textit{"Ele estudou de manhã; ela, à tarde."} (omitiu-se
"estudou" na segunda oração --- também chamado de zeugma quando a omissão é
de um termo já citado).
\end{itemize}

O que retoma pode ser um \textbf{pronome} (ele, o qual, isso), um
\textbf{sinônimo} ou \textbf{hiperônimo} (repetir "o documento" chamando de
"o papel" ou de "o texto"), ou uma \textbf{elipse}.
\end{conceito}

\subsection{Conectores e o valor semântico que carregam}

A reescrita/interpretação cobra menos o nome da conjunção e mais \textbf{que
relação lógica} ela expressa --- porque a mesma relação pode ser expressa por
conectores de classes gramaticais diferentes.

\begin{center}
\begin{tabular}{@{}p{2.6cm}p{9.4cm}@{}}
\toprule
\textbf{Relação} & \textbf{Conectores que a expressam} \\
\midrule
Adição & e, nem, também, além disso, ainda \\
Oposição/contraste & mas, porém, contudo, todavia, entretanto, no entanto,
embora, ainda que, apesar de \\
Causa & porque, já que, uma vez que, visto que, por causa de, devido a \\
Consequência & portanto, logo, por isso, de modo que, de forma que \\
Condição & se, caso, contanto que, desde que, a menos que \\
Finalidade & para que, a fim de que, com o objetivo de \\
Tempo & quando, enquanto, assim que, logo que, antes que, depois que \\
Comparação & como, tal qual, assim como, mais/menos... do que \\
Conformidade & conforme, segundo, de acordo com, consoante \\
\bottomrule
\end{tabular}
\end{center}

\begin{cuidado}[O mesmo conector muda de valor com o contexto]
"Como" pode ser causal ("Como chovia, não saiu" = "porque chovia"),
comparativo ("Ele é como o pai") ou conformativo ("Como diz o ditado...").
"Que" pode ser conjunção integrante (introduz substantiva), pronome relativo
(introduz adjetiva) ou até conjunção consecutiva ("tão cansado que
dormiu"). A prova nunca pergunta só "que classe é esta palavra" --- pergunta
"que \textbf{relação}/\textbf{função} ela introduz \textbf{neste}
contexto", e a resposta muda frase a frase.
\end{cuidado}

\subsection{Correlação/sequência de tempos verbais}
\label{sec:correlacao}

\begin{conceito}[O tempo da oração principal governa o da subordinada]
Quando duas orações do período estão relacionadas no tempo, o verbo da
subordinada tem que ser compatível com o tempo do verbo da principal --- é a
mesma lógica do "sequence of tenses" do inglês, só que menos rígida em
português. Padrões mais cobrados:

\begin{itemize}
\item Principal no \textbf{presente} $\to$ subordinada pode ir para
qualquer tempo compatível com o sentido: \textit{"Acho que ele
\textbf{chegou}/\textbf{chega}/\textbf{vai chegar}."}
\item Principal no \textbf{pretérito} $\to$ a subordinada tende a
"puxar" para o pretérito também (discurso indireto clássico):
\textit{"Ele disse que \textbf{viria}"} (não "que virá") --- futuro do
pretérito no lugar do futuro do presente, porque a fala foi relatada no
passado.
\item \textbf{Condicional}: "se" + \textbf{pretérito imperfeito do
subjuntivo} na condição pede \textbf{futuro do pretérito} (não futuro do
presente) na consequência: \textit{"Se eu \textbf{tivesse} tempo, eu
\textbf{estudaria}"} (não "eu estudarei"). É um dos erros mais comuns de
concordância verbal em prova de reescrita.
\end{itemize}
\end{conceito}

\section{Classes de palavras}

\begin{conceito}[As dez classes, divididas em variáveis e invariáveis]
\textbf{Variáveis} (flexionam em gênero, número, e às vezes grau/pessoa/modo/tempo):
substantivo, artigo, adjetivo, numeral, pronome, verbo.
\textbf{Invariáveis} (forma fixa): advérbio, preposição, conjunção,
interjeição.

O ponto que a prova explora não é decorar a lista, é reconhecer que
\textbf{a mesma palavra muda de classe conforme a função na frase}
(derivação imprópria / hibridismo funcional):
\begin{itemize}
\item "Amanhã" $\to$ substantivo ("o amanhã é incerto") ou advérbio
("chego amanhã").
\item "Que" $\to$ pronome relativo, conjunção integrante, conjunção
consecutiva, ou interjeição de espanto ("Que confusão!").
\item "Mal" $\to$ advérbio ("ele foi mal") ou substantivo ("um mal
necessário") --- não confundir com "mau" (adjetivo, antônimo de "bom").
\end{itemize}
\end{conceito}

\begin{regrapratica}
Para classificar uma palavra ambígua, \textbf{substitua-a mantendo a
função}: se dá para trocar por "isso"/"aquilo" e ainda ocupar posição de
sujeito/objeto, é substantivo; se modifica um verbo/adjetivo/outro advérbio
sem flexionar, é advérbio. Nunca classifique a palavra isolada --- classifique
a palavra \textbf{na frase dada}.
\end{regrapratica}

\section{Sintaxe do período: coordenação e subordinação}

\subsection{Oração e período}

\begin{conceito}[Contando orações]
\textbf{Oração} é todo enunciado organizado em torno de um núcleo verbal
(verbo ou locução verbal). \textbf{Período simples} tem uma oração só;
\textbf{período composto} tem duas ou mais. Truque de contagem: conte os
verbos/locuções verbais do período --- cada núcleo verbal é, em regra, uma
oração (locução verbal conta como uma só).
\end{conceito}

\begin{exemplo}
\textit{"O candidato estudou muito porque a prova é difícil."} $\to$ dois
núcleos verbais ("estudou", "é") $\to$ período composto de duas orações.
\end{exemplo}

O período composto se organiza por dois mecanismos: \textbf{coordenação}
(orações independentes entre si) e \textbf{subordinação} (uma oração é termo
sintático da outra). Os dois podem coexistir no mesmo período (período
misto).

\subsection{Coordenação}

\begin{conceito}[Orações coordenadas]
São \textbf{sintaticamente independentes}: nenhuma delas exerce função de
termo (sujeito, objeto etc.) dentro da outra --- cada uma "sobrevive
sozinha" gramaticalmente. Podem vir \textbf{sem} conectivo
(\textbf{assindéticas}, separadas por vírgula/ponto e vírgula: \textit{"Cheguei,
vi, venci"}) ou \textbf{com} conjunção coordenativa (\textbf{sindéticas}),
que se dividem em cinco tipos:

\begin{center}
\begin{tabular}{@{}p{2.6cm}p{4.6cm}p{4.8cm}@{}}
\toprule
\textbf{Tipo} & \textbf{Conjunções típicas} & \textbf{Ideia} \\
\midrule
Aditiva & e, nem, mas também & soma \\
Adversativa & mas, porém, contudo, todavia, entretanto, no entanto & contraste/oposição \\
Alternativa & ou, ou...ou, ora...ora, seja...seja & escolha/alternância \\
Conclusiva & logo, portanto, por isso, pois (depois do verbo) & conclusão \\
Explicativa & pois (antes do verbo), porque, que & justificativa \\
\bottomrule
\end{tabular}
\end{center}
\end{conceito}

\begin{cuidado}[A conjunção "pois" muda de classe com a posição]
\textit{"Estude, pois a prova é difícil."} --- "pois" logo depois da vírgula,
\textbf{antes} do verbo $\to$ \textbf{explicativa} (equivale a "porque").
\textit{"A prova é difícil; estude, pois."} / \textit{"Estude; a prova, pois,
é difícil."} --- "pois" deslocado, intercalado \textbf{no meio ou no fim} da
oração $\to$ \textbf{conclusiva} (equivale a "portanto"). Regra prática:
pois colado após a vírgula = explicativa; pois "solto" no meio da frase =
conclusivo.
\end{cuidado}

\subsection{Subordinação}

\begin{conceito}[Orações subordinadas]
Uma oração \textbf{depende sintaticamente} da outra: ela ocupa, sozinha, a
função que uma \textbf{palavra} ocuparia (sujeito, objeto, adjunto...). A
pergunta-chave para classificar é sempre: \textbf{"que classe gramatical essa
oração inteira está imitando?"} --- substantivo, adjetivo ou advérbio. Daí os
três grandes grupos: \textbf{substantivas}, \textbf{adjetivas},
\textbf{adverbiais}.
\end{conceito}

\subsubsection{Substantivas}
\label{sec:substantivas}

\begin{conceito}[A oração no lugar de um substantivo]
Iniciadas por conjunção \textbf{integrante} ("que" ou "se" --- não confundir
com pronome relativo). Teste rápido: troque a oração inteira por
\textbf{"isso"}; se a frase continuar de pé, é substantiva, e a posição de
"isso" na frase revela a função:

\begin{center}
\begin{tabular}{@{}p{2.6cm}p{2.8cm}p{6.6cm}@{}}
\toprule
\textbf{Subtipo} & \textbf{Função} & \textbf{Exemplo} \\
\midrule
Subjetiva & sujeito & \textit{É necessário \textbf{que você estude}.} ("Isso é necessário") \\
Objetiva direta & objeto direto & \textit{Espero \textbf{que você passe}.} ("Espero isso") \\
Objetiva indireta & objeto indireto & \textit{Preciso \textbf{de que você me ajude}.} ("Preciso disso") \\
Completiva nominal & completa um nome & \textit{Tenho certeza \textbf{de que ele passou}.} ("Tenho certeza disso") \\
Predicativa & predicativo do sujeito & \textit{O problema é \textbf{que faltou tempo}.} ("O problema é isso") \\
Apositiva & aposto & \textit{Só quero uma coisa: \textbf{que você estude}.} \\
\bottomrule
\end{tabular}
\end{center}
\end{conceito}

\begin{cuidado}[Objetiva indireta x completiva nominal]
As duas têm preposição, e é fácil trocar uma pela outra. O critério: se a
preposição é \textbf{exigida pelo verbo} (regência verbal), é objetiva
indireta; se é exigida por um \textbf{nome abstrato} da oração principal
(substantivo/adjetivo como "certeza", "medo", "necessidade", "capaz"), é
completiva nominal. \textit{"Preciso \underline{de} que você ajude"} (verbo
"precisar" rege "de") é objetiva indireta; \textit{"Tenho certeza
\underline{de} que ele passou"} ("certeza de" --- o nome "certeza" rege "de")
é completiva nominal.
\end{cuidado}

\subsubsection{Adjetivas}

\begin{conceito}[A oração no lugar de um adjetivo]
Sempre iniciadas por \textbf{pronome relativo} (que, quem, o qual, cujo,
onde, quando referidos a um antecedente), caracterizam um substantivo/pronome
da oração principal (o \textbf{antecedente}). Dividem-se em:

\begin{itemize}
\item \textbf{Restritivas} (\textbf{sem} vírgula): delimitam o universo do
antecedente --- retirar a oração \textbf{muda o sentido} da frase.
\textit{"Os candidatos que estudaram passaram."} (só os que estudaram, não
todos).
\item \textbf{Explicativas} (\textbf{com} vírgula): acrescentam uma
informação extra sobre \textbf{todo} o antecedente --- retirar a oração
\textbf{não muda} o sentido essencial. \textit{"Os candidatos, que
estudaram muito, passaram."} (todos estudaram; é um comentário à parte).
\end{itemize}
\end{conceito}

\begin{cuidado}[A vírgula muda o sentido da frase inteira]
Este é o ponto mais cobrado do capítulo inteiro de sintaxe: a mesma
sequência de palavras, com e sem vírgula, afirma coisas diferentes.
\textit{"Os alunos que estudaram passaram"} não garante que todos os alunos
passaram (só os que estudaram); \textit{"Os alunos, que estudaram,
passaram"} afirma que \textbf{todos} os alunos estudaram e \textbf{todos}
passaram. Ao comparar duas versões de uma frase numa questão de reescrita,
a vírgula antes do "que" é o primeiro lugar a olhar.
\end{cuidado}

\begin{regrapratica}[Cujo]
"Cujo" é sempre pronome relativo (introduz adjetiva), concorda em
gênero/número com a coisa \textbf{possuída} (não com o possuidor), e nunca
vem seguido de artigo. \textit{"O candidato cujo caderno sumiu"} (não "cujo
o caderno").
\end{regrapratica}

\subsubsection{Adverbiais}

\begin{conceito}[A oração no lugar de um advérbio]
Fazem o papel de adjunto adverbial (circunstância) da oração principal. Nove
tipos, decorados pela conjunção-gatilho:

\begin{center}
\begin{tabular}{@{}p{2.6cm}p{4.8cm}p{4.6cm}@{}}
\toprule
\textbf{Tipo} & \textbf{Conjunções típicas} & \textbf{Ideia} \\
\midrule
Causal & porque, já que, uma vez que, visto que & causa \\
Consecutiva & que (após "tão, tal, tanto") & consequência \\
Condicional & se, caso, contanto que, desde que & condição \\
Concessiva & embora, ainda que, mesmo que & concessão \\
Comparativa & como, (do) que & comparação \\
Conformativa & conforme, segundo, como & conformidade \\
Final & para que, a fim de que & finalidade \\
Temporal & quando, enquanto, assim que, logo que & tempo \\
Proporcional & à medida que, à proporção que & proporção \\
\bottomrule
\end{tabular}
\end{center}
\end{conceito}

\begin{cuidado}[Causal x consecutiva]
\textit{"Estudou tanto que passou."} \textbf{Não} é causal --- é
\textbf{consecutiva}: a conjunção correlata é "tanto...que", e a ideia é
consequência (o resultado de estudar tanto foi passar), não a causa de algo.
\end{cuidado}

\begin{cuidado}[Concessiva não é "adversativa disfarçada"]
As duas contrapõem ideias, mas por mecanismos sintáticos diferentes.
\textbf{Adversativa} é \textbf{coordenação}: duas orações independentes
contrapostas (\textit{"Estudou, mas não passou"}). \textbf{Concessiva} é
\textbf{subordinação}: uma oração se subordina admitindo um obstáculo que
não impede a outra (\textit{"Embora tenha estudado, não passou"}). Trocar
"mas" por "embora" (ou vice-versa) muda a estrutura sintática do período
inteiro, não só o conectivo.
\end{cuidado}

\subsection{Orações reduzidas}

\begin{conceito}[Sem conectivo, sem verbo flexionado]
Quando a subordinada não tem conectivo nem verbo flexionado (infinitivo,
gerúndio ou particípio), ela é \textbf{reduzida}. A versão com conectivo e
verbo flexionado é a \textbf{desenvolvida}.

\begin{center}
\begin{tabular}{@{}p{2.6cm}p{4.2cm}p{5.2cm}@{}}
\toprule
\textbf{Forma nominal} & \textbf{Reduzida} & \textbf{Desenvolvida equivalente} \\
\midrule
Infinitivo & \textit{Ao chegar, avisou.} & \textit{Quando chegou, avisou.} (valor temporal) \\
Gerúndio & \textit{Estudando, você passa.} & \textit{Se você estudar, você passa.} (valor condicional) \\
Particípio & \textit{Terminada a prova, saiu.} & \textit{Quando a prova terminou, saiu.} (valor temporal) \\
\bottomrule
\end{tabular}
\end{center}
\end{conceito}

\begin{regrapratica}
Para classificar uma reduzida, \textbf{desenvolva-a mentalmente} (acrescente
o conectivo que faz sentido) e classifique pelo \textbf{valor semântico}
resultante (causal, temporal, condicional...) --- a forma nominal
(infinitivo/gerúndio/particípio) é só a "casca"; o sentido é que define o
subtipo, porque não há conjunção explícita ali para entregar a resposta.
\end{regrapratica}

\subsection{Roteiro de classificação}

\begin{regrapratica}[Ordem de decisão para classificar um período composto]
\begin{enumerate}
\item Conte os núcleos verbais $\to$ quantas orações tem o período.
\item Cada oração (a partir da segunda) é independente ou depende de outra
como se fosse um termo dela?
\begin{itemize}
\item Independente $\to$ \textbf{coordenada}: identifique a conjunção,
classifique (aditiva/adversativa/alternativa/conclusiva/explicativa) ou
marque assindética.
\item Depende (funciona como termo da outra) $\to$ \textbf{subordinada}:
pergunte que classe gramatical ela imita.
\begin{itemize}
\item Substantivo $\to$ substantiva $\to$ identifique a função sintática
(troque por "isso") $\to$ subtipo.
\item Adjetivo (tem antecedente + pronome relativo) $\to$ adjetiva $\to$ tem
vírgula? restritiva/explicativa.
\item Advérbio (circunstância) $\to$ adverbial $\to$ identifique a
conjunção-gatilho $\to$ subtipo.
\end{itemize}
\end{itemize}
\item Verbo no infinitivo/gerúndio/particípio sem conectivo? $\to$
\textbf{reduzida}, classifique pelo sentido (desenvolva mentalmente).
\end{enumerate}
\end{regrapratica}

\section{Pontuação}

\begin{conceito}[A vírgula tem uma lista curta e fechada de usos válidos]
Cada vírgula "correta" cabe em uma destas categorias --- se não cabe em
nenhuma, está errada:

\begin{itemize}
\item \textbf{Enumeração}: separa itens de uma série (\textit{"Trouxe
caneta, lápis, borracha e régua."} --- repare que não há vírgula antes do
último item ligado por "e", salvo se os itens tiverem sujeitos diferentes).
\item \textbf{Aposto}: explica/identifica um termo, isolado por vírgulas
(\textit{"Machado de Assis, o maior escritor brasileiro, nasceu no Rio."}).
\item \textbf{Vocativo}: chama alguém (\textit{"Lucas, venha aqui."}).
\item \textbf{Oração intercalada}: interrompe a frase principal
(\textit{"O réu, segundo o juiz, mentiu."}).
\item \textbf{Adjunto adverbial deslocado}: quando o adjunto sai do lugar
padrão (fim da frase) e vai para o início/meio (\textit{"Ontem, o sistema
caiu."} --- vírgula opcional se curto, recomendada se longo).
\item \textbf{Oração subordinada adverbial anteposta à principal}
(\textit{"Se chover, cancelamos o evento."} --- a vírgula aqui é
\textbf{obrigatória}; se a ordem inverter, \textit{"Cancelamos o evento se
chover"}, some).
\item \textbf{Oração adjetiva explicativa} (ver seção anterior).
\item \textbf{Elipse de verbo} (\textit{"Ele estudou de manhã; ela, à
tarde."}).
\end{itemize}
\end{conceito}

\begin{cuidado}[Vírgula proibida]
Nunca se separa por vírgula: sujeito do predicado (\textit{"O candidato,
passou"} está errado), verbo do complemento (\textit{"Ele precisa, de
ajuda"} está errado), salvo quando entre eles se intercala aposto ou oração
intercalada (aí a vírgula existe por causa do intercalado, não "entre sujeito
e verbo").
\end{cuidado}

\begin{conceito}[Ponto e vírgula e dois-pontos]
\textbf{Ponto e vírgula}: separa orações coordenadas quando pelo menos uma já
tem vírgula interna (evita ambiguidade), ou itens de uma enumeração longa
com sub-itens. \textbf{Dois-pontos}: anuncia uma enumeração, uma citação, uma
explicação ou um aposto que resume o que veio antes (\textit{"Só quero uma
coisa: que você estude."}).
\end{conceito}

\section{Concordância verbal}

\subsection{Regra geral e sujeito composto}

\begin{conceito}[O verbo concorda com o sujeito]
Regra base: o verbo flexiona em número e pessoa para concordar com o núcleo
do sujeito. Com \textbf{sujeito composto} anteposto ao verbo, o verbo vai
para o \textbf{plural} (\textit{"João e Maria \textbf{chegaram}."}). Com
sujeito composto \textbf{posposto} ao verbo, admite-se concordância no plural
\textbf{ou} com o núcleo mais próximo (\textit{"\textbf{Chegaram}
João e Maria"} ou, mais raro em prova, \textit{"Chegou João e Maria"}).
\end{conceito}

\subsection{Verbos impessoais}

\begin{conceito}[Verbos sem sujeito ficam travados na 3ª pessoa do singular]
Um pequeno grupo de verbos, em certos sentidos, \textbf{não tem sujeito}
--- e por isso nunca vão para o plural, não importa o que venha depois deles:

\begin{center}
\begin{tabular}{@{}p{2.6cm}p{3.6cm}p{5.2cm}@{}}
\toprule
\textbf{Verbo} & \textbf{Quando é impessoal} & \textbf{Certo (e o erro comum)} \\
\midrule
\textbf{fazer} & tempo decorrido & \textit{Faz} dois anos --- nunca "Fazem dois anos" \\
\textbf{haver} & sentido de "existir" ou tempo & \textit{Havia} muitas falhas --- nunca "Haviam muitas falhas" \\
\textbf{haver} + auxiliar & o auxiliar também trava & \textit{Deve haver} relatórios --- nunca "Devem haver" \\
\textbf{existir} & \textbf{não} é impessoal, tem sujeito & \textit{Existem} muitos processos --- nunca "Existe muitos processos" \\
\bottomrule
\end{tabular}
\end{center}

O contraste que a banca explora: "haver" (no sentido de existir) e "existir"
significam quase a mesma coisa, mas um é impessoal (não varia) e o outro tem
sujeito (varia normalmente).
\end{conceito}

\subsection{Sujeito coletivo e sujeito oracional}

\begin{conceito}
\textbf{Sujeito coletivo} (multidão, maioria, grupo, equipe) concorda no
\textbf{singular} quando não especificado (\textit{"A maioria
\textbf{chegou}."}), mas admite plural quando o coletivo vem seguido de um
complemento no plural (\textit{"A maioria dos candidatos
\textbf{chegou}/\textbf{chegaram}"} --- as duas são aceitas; a FGV costuma
cobrar a forma no singular como a "mais formal/tradicional").
\textbf{Sujeito oracional}: quando o sujeito é uma oração subordinada
substantiva subjetiva inteira, o verbo da principal fica sempre no
\textbf{singular} (\textit{"\textbf{É} preciso que todos estudem."} --- não
"são preciso").
\end{conceito}

\section{Concordância nominal}

\begin{conceito}[Adjetivo concorda com o substantivo a que se refere]
Regra geral simples; os casos que a prova cobra são os que \textbf{parecem}
exceção mas seguem lógica própria:

\begin{itemize}
\item \textbf{"Anexo", "incluso", "obrigado", "próprio", "mesmo", "quite"}:
são \textbf{adjetivos} e concordam normalmente (\textit{"Seguem
\textbf{anexas} as fotos."}, \textit{"Ela mesma resolveu."}, \textit{"Ele
está \textbf{obrigado}, ela está \textbf{obrigada}."}) --- o erro comum é
tratá-los como advérbios invariáveis.
\item \textbf{"Meio" advérbio (= "um pouco")} é \textbf{invariável}
(\textit{"Ela está \textbf{meio} confusa."} --- não "meia confusa"); "meio"
\textbf{numeral} (= metade) concorda (\textit{"Comeu \textbf{meia}
maçã."}).
\item \textbf{"É proibido/necessário/preciso" + substantivo}: fica
\textbf{invariável} (masculino singular) quando o substantivo \textbf{não}
tem determinante (artigo/pronome) antes (\textit{"\textbf{É proibido}
entrada."}); \textbf{concorda} quando o substantivo vem com determinante
(\textit{"\textbf{É proibida} \underline{a} entrada."}).
\item \textbf{"Bastante"}: como advérbio (modifica verbo/adjetivo) é
invariável (\textit{"Elas estudaram \textbf{bastante}."}); como
adjetivo/pronome (modifica substantivo, = "muitos/muitas") concorda
(\textit{"\textbf{Bastantes} candidatos desistiram."}).
\end{itemize}
\end{conceito}

\begin{cuidado}[O padrão por trás de "é proibido"/"é bastante"]
Não são regras soltas: a lógica comum é \textbf{"a palavra concorda quando
funciona como adjetivo do substantivo; fica fixa quando funciona como
advérbio/predicativo genérico sem se prender a um substantivo
determinado"}. Uma vez internalizado esse padrão, "meio", "bastante" e "é
proibido" deixam de ser três regras decoradas e viram uma regra só aplicada
três vezes.
\end{cuidado}

\section{Regência verbal e nominal}

\begin{conceito}[Regência é sobre a preposição que o verbo/nome exige]
\textbf{Regência verbal}: a preposição (ou ausência dela) que cada verbo
exige para introduzir seu complemento --- e o mesmo verbo pode mudar de
sentido conforme a regência:

\begin{center}
\begin{tabular}{@{}p{2.4cm}p{4.6cm}p{4.4cm}@{}}
\toprule
\textbf{Verbo} & \textbf{Regência e sentido} & \textbf{Exemplo} \\
\midrule
assistir & "a" = ver/presenciar (transitivo indireto) & Assisti \textbf{ao} filme. \\
assistir & sem prep.\ = ajudar, dar assistência (transitivo direto) & O médico assistiu o paciente. \\
visar & "a" = ter como objetivo (formal) & O projeto visa \textbf{a} reduzir custos. \\
visar & sem prep.\ = mirar, rubricar & Visou o documento; visou o alvo. \\
obedecer / desobedecer & sempre "a" (transitivo indireto) & Obedeça \textbf{às} normas. \\
esquecer & sem prep.\ (transitivo direto) OU pronominal "esquecer-se de" & Esqueci o nome. / Esqueci-me \textbf{do} nome. \\
lembrar & sem prep.\ (transitivo direto) OU pronominal "lembrar-se de" & Lembrei o nome. / Lembrei-me \textbf{do} nome. \\
implicar & sem prep.\ = acarretar (transitivo direto, sentido de "causar") & O atraso implicou \textbf{a} perda do prazo. \\
preferir & "a" (nunca "do que") na norma culta & Prefiro café \textbf{a} chá. \\
\bottomrule
\end{tabular}
\end{center}

\textbf{Regência nominal}: substantivos e adjetivos abstratos também exigem
preposição fixa (\textit{"obediente \textbf{a}"}, \textit{"apto
\textbf{para}"}, \textit{"cheio \textbf{de}"}, \textit{"certeza
\textbf{de}"}) --- é a mesma preposição que aparece na oração completiva
nominal (ver Seção~\ref{sec:substantivas} deste capítulo).
\end{conceito}

\begin{cuidado}[Regência com duplo regime é área de controvérsia gramatical]
Alguns pares de regência têm \textbf{mais de uma forma aceita} pela norma
culta (ex.: "namorar" pode ou não levar preposição conforme a gramática
consultada; "chegar/ir a/para" varia por região e por tradição/permanência
pretendida). Nesses casos não existe "a" resposta absoluta --- a banca, ao
cobrar, costuma usar o caso \textbf{mais consolidado} (assistir, obedecer,
visar, preferir da tabela acima), não os limítrofes. Se um caso limítrofe
aparecer, desconfie de alternativas que tratam a variação como erro
categórico.
\end{cuidado}

\section{Crase}

\begin{conceito}[Crase é a fusão da preposição "a" com o "a" artigo/pronome]
Crase só existe quando \textbf{duas condições se encontram ao mesmo tempo}:
(1) um termo antes que \textbf{exige a preposição "a"} (regência verbal ou
nominal) e (2) um termo depois que \textbf{admite artigo feminino "a(s)"}.
Faltando qualquer uma das duas, não há crase.

\begin{itemize}
\item \textbf{Obrigatória}: locuções adverbiais femininas de instrumento/modo/causa
("à mão", "às pressas", "à força"); locuções conjuntivas/prepositivas
femininas ("à medida que", "à proporção que"); antes de "moda"/"maneira"
subentendida ("sapato à Luís XV" = "à moda de"); horas determinadas quando
o contexto pede "a" ("Cheguei às 8h").
\item \textbf{Proibida}: antes de verbo (\textit{"Começou \textbf{a}
estudar"}, nunca "à estudar"); antes de palavra masculina (\textit{"Andar
\textbf{a} cavalo"}); antes de pronome que não admite artigo (\textit{"Entreguei
\textbf{a} ela"}, nunca "à ela" --- pronomes pessoais não levam artigo,
salvo exceções específicas); entre palavras repetidas (\textit{"cara a
cara", "dia a dia"}); antes de "casa" (sem especificador) e "terra" (sentido
de chão firme, oposto a bordo) no sentido genérico.
\item \textbf{Facultativa}: antes de pronome possessivo feminino
(\textit{"Entreguei a/à sua filha"}); antes de nome próprio de mulher
(\textit{"Entreguei a/à Maria"}); depois da preposição "até"
(\textit{"Foi até a/à porta"}).
\end{itemize}
\end{conceito}

\begin{regrapratica}[Teste da substituição por palavra masculina]
Troque o termo depois do "a" por um equivalente \textbf{masculino}. Se o
masculino pede "ao", o feminino leva crase; se o masculino pede só "o" (sem
preposição) ou "a" (sem artigo), não há crase. \textit{"Foi à escola"}
$\to$ masculino: "Foi \textbf{ao} colégio" $\to$ crase confirmada.
\textit{"Foi a pé"} $\to$ masculino: "Foi \textbf{a} cavalo" (não "ao
cavalo") $\to$ sem crase.
\end{regrapratica}

\begin{cuidado}[Crase em caso limítrofe também tem controvérsia gramatical]
Casos como crase antes de nomes de lugares que admitem artigo de forma
variável ("Vou a Bahia" x "Vou à Bahia", dependendo se o falante entende
"Bahia" como admitindo artigo) ou antes de pronomes possessivos são
\textbf{genuinamente controversos} entre gramáticos — não é uma lacuna do
material, é a língua tendo duplo regime ali. A banca, quando quer resposta
inequívoca, usa os casos obrigatórios/proibidos da lista acima, não os
facultativos.
\end{cuidado}

\section{Colocação pronominal}

\begin{conceito}[Próclise, mesóclise e ênclise]
Trata da posição do pronome oblíquo átono (me, te, se, o, a, lhe, nos, vos,
lhes) em relação ao verbo:

\begin{itemize}
\item \textbf{Próclise} (pronome \textbf{antes} do verbo): puxada por uma
\textbf{palavra atrativa} antes do verbo --- palavra de sentido negativo
("Não \textbf{me} diga"), pronome relativo ("o homem que \textbf{me}
ajudou"), conjunção subordinativa ("quando \textbf{se} formou"), advérbio
sem vírgula depois ("Aqui \textbf{se} vive bem"), pronome indefinido ou
interrogativo ("Ninguém \textbf{me} avisou", "Quem \textbf{te} contou?").
\item \textbf{Mesóclise} (pronome \textbf{no meio} do verbo): só ocorre com
verbo no \textbf{futuro do presente} ou \textbf{futuro do pretérito}, e
apenas quando não há palavra atrativa antes (\textit{"Dir-\textbf{lhe}-ei a
verdade"}; se houvesse atrativa antes, viraria próclise: "Não \textbf{lhe}
direi").
\item \textbf{Ênclise} (pronome \textbf{depois} do verbo): a posição
"neutra"/padrão quando não há atrativa antes e o verbo não está no
futuro --- é a posição de \textbf{início de frase} (\textit{"\textbf{Diga}-me
a verdade"}) e a recomendada na norma culta escrita quando não há motivo
para próclise.
\end{itemize}
\end{conceito}

\begin{cuidado}[Início de frase nunca leva próclise na norma culta escrita]
\textit{"Me diga a verdade"} é falado corrente no Brasil, mas na norma culta
escrita (a que a prova cobra) frase iniciada por verbo exige ênclise:
\textbf{"Diga-me a verdade."} É um dos pontos onde o português informal
brasileiro mais diverge da norma cobrada em concurso.
\end{cuidado}

\section{Reescrita e significação}

\begin{conceito}[O que muda e o que não pode mudar numa reescrita]
Questão de reescrita pede para reescrever um trecho \textbf{mudando a forma
sem mudar o sentido} (ou, menos comum, mudando intencionalmente o sentido, e
aí a pergunta avisa). O que se testa:

\begin{itemize}
\item \textbf{Sinonímia}: troca de palavra por um sinônimo que preserva o
sentido \textbf{no contexto} (nem todo sinônimo de dicionário serve em toda
frase --- "grande" pode ser "amplo" em um contexto e "importante" em outro).
\item \textbf{Denotação x conotação}: sentido \textbf{literal/dicionário}
(denotação) x sentido \textbf{figurado/contextual} (conotação). Reescrever
trocando um sentido conotativo por uma paráfrase literal do mesmo sentido
(não pelo sentido literal da palavra) é o padrão mais cobrado.
\item \textbf{Ambiguidade}: mais de uma leitura sintática/semântica possível
para a mesma frase --- reescrita "corrige" ambiguidade reordenando termos ou
trocando um pronome vago por um termo específico (\textit{"João disse a
Pedro que \textbf{ele} passou"} --- "ele" é ambíguo; a reescrita sem
ambiguidade precisa nomear quem passou).
\item \textbf{Voz ativa x passiva (apassivação)}: mesma informação,
estrutura sintática diferente (\textit{"O réu cometeu o crime"} $\leftrightarrow$
\textit{"O crime foi cometido pelo réu"}) --- o sujeito da ativa vira agente
da passiva, o objeto direto da ativa vira sujeito da passiva.
\item \textbf{Discurso direto x indireto}: o direto reproduz literalmente a
fala (aspas/travessão, verbo de elocução: "disse:", "perguntou:"); o
indireto relata a fala dentro de uma oração subordinada, com os ajustes de
pessoa e tempo verbal da Seção~\ref{sec:correlacao} (\textit{"Ele disse: 'Eu
vou'"} $\to$ \textit{"Ele disse que \textbf{iria}"}).
\end{itemize}
\end{conceito}

\begin{exemplo}[Reescrita preservando sentido]
Original: \textit{"Embora tenha estudado bastante, ele não passou."}
Reescritas equivalentes: \textit{"Ele estudou bastante, mas não passou."}
(concessiva $\to$ adversativa, mudou de subordinação para coordenação, mesmo
sentido); \textit{"Apesar de ter estudado bastante, ele não passou."}
(troca de conjunção concessiva por locução prepositiva + infinitivo,
reduzida).
\end{exemplo}

\begin{regrapratica}
Para validar uma reescrita, pergunte: \textbf{quem} fez \textbf{o quê}
\textbf{a quem}, mudou? Se as relações de sujeito/objeto/causa/consequência
continuam as mesmas e só a \textbf{forma} (ordem das palavras, voz, classe
gramatical do conectivo) mudou, a reescrita é válida. Se alguma relação
lógica se inverteu ou sumiu, não é reescrita equivalente --- é outra frase.
\end{regrapratica}
